\par\bigskip\Needspace{7\baselineskip}
\pdfbookmark[1]{Handbuch der Schiffahrtskunde. 1850.}{schiffahrtskunde-1850}
\begin{center}Handbuch der Schiffahrtskunde von C. \textsc{Rümker}. 1850. Seite 76.\par\medskip
\textit{Lösung einer geometrischen Aufgabe.}\end{center}
\hspace*{\parindent}An drei Punkten \((1),(2),(3)\), die auf einer Geraden \(\mathrm{(I)}\) in bekannten Abständen voneinander liegen, \(A\) von \((1)\) nach \((2)\) und \(B\) von \((2)\) nach \((3)\), wurden die Winkel \(\vartheta,\vartheta',\vartheta''\) zwischen zwei anderen Punkten \((4),(5)\) gemessen, deren gegenseitiger Abstand \(=2c\) ebenfalls bekannt ist. Gesucht ist die Lage der ersten drei Punkte relativ zu den beiden letzten. Damit nichts unbestimmt bleibt, setze ich voraus, dass alle drei Winkel von \((4)\) nach \((5)\) im selben Drehsinn zunehmend gemessen wurden und dass auf der Geraden \(\mathrm{(I)}\) die Abstände in einer festgelegten Richtung positiv gezählt werden. Würde man also aus irgendeinem Grund den Punkt zwischen den beiden anderen anders als mit \((2)\) bezeichnen, hätten \(A\) und \(B\) verschiedene Vorzeichen. Außerdem soll \(c\) positiv sein.
